\chapter{Mathematical toolkit}
\label{ch:math}

This chapter collects every piece of mathematics used later. If a symbol in a later chapter looks
unfamiliar, it is defined here.

\section{Scalars, vectors, fields}

A \textbf{scalar} is a single number (temperature $T$, density $\rho$). A \textbf{vector} has
magnitude and direction; in 3D Cartesian coordinates $\vu=(u,v,w)$. A \textbf{field} is a quantity
defined at every point and time: $T(\vx,t)$ is a scalar field, $\vu(\vx,t)$ a vector field.

The \textbf{dot product} $\bm a\cdot\bm b = a_x b_x + a_y b_y + a_z b_z = |\bm a||\bm b|\cos\theta$
measures how much two vectors point the same way. With a unit normal $\vn$ to a surface,
$\vu\cdot\vn$ is the velocity component \emph{through} that surface.

\paragraph{Index notation.} Writing $u_i$ ($i=1,2,3$) for the components and using the
\emph{summation convention} (a repeated index is summed), $\bm a\cdot \bm b = a_i b_i$. We use it
only in the momentum equation.

\section{The gradient: which way is uphill}

For a scalar field $T(x,y,z)$,
\begin{equation}
  \nabla T = \left(\pd{T}{x},\pd{T}{y},\pd{T}{z}\right).
\end{equation}
$\nabla T$ points in the direction of steepest increase of $T$ and its magnitude is the rate of
increase per metre. The change of $T$ over a small displacement $\dd\vx$ is
$\dd T = \nabla T\cdot \dd \vx$.

\begin{intuition}
Heat flows from hot to cold, i.e.\ \emph{down} the gradient. That is why Fourier's law
(\chapref{ch:transport}) is $\vq = -\lambda\nabla T$, with a minus sign. Species diffuse down their
concentration gradient in the same way. In a flame, the gradients are enormous: temperature rises by
$\sim\SI{2000}{\kelvin}$ over $\sim\SI{0.5}{\milli\metre}$, i.e.\ $|\nabla T|\sim\SI{4e6}{\kelvin\per\metre}$.
\end{intuition}

\section{The divergence: net outflow per unit volume}
\label{sec:divergence}

Let $\bm F$ be a \emph{flux}: an amount of something crossing a unit area per unit time (e.g.\
heat flux in \si{\watt\per\metre\squared}, mass flux $\rho\vu$ in \si{\kilo\gram\per\metre\squared\per\second}).

\begin{figure}[H]
\centering
\begin{tikzpicture}[scale=1.2]
  \draw[thick,fill=gray!8] (0,0) rectangle (3,2);
  \draw[->,very thick,cold] (-1.3,1) -- (0,1) node[midway,above,lbl] {$F_x(x)$};
  \draw[->,very thick,cold] (3,1) -- (4.6,1) node[midway,above,lbl] {$F_x(x+\Delta x)$};
  \draw[<->] (0,-0.3) -- (3,-0.3) node[midway,below,lbl] {$\Delta x$};
  \draw[<->] (-0.3,0) -- (-0.3,2) node[midway,left,lbl] {$\Delta y$};
  \node[lbl] at (1.5,1) {volume $\Delta V$};
  \draw[->,very thick,cold] (1.5,2) -- (1.5,2.8) node[right,lbl] {$F_y(y+\Delta y)$};
  \draw[->,very thick,cold] (1.5,-1.0) -- (1.5,0) node[pos=0.1,right,lbl] {$F_y(y)$};
\end{tikzpicture}
\caption{A small box. What leaves through the right face minus what enters through the left face,
plus the same in $y$ and $z$, is the net outflow.}
\label{fig:divbox}
\end{figure}

\begin{derivation}[divergence from a box]
Net outflow in $x$ per unit time through faces of area $\Delta y\Delta z$:
\[
 \big[F_x(x+\Delta x)-F_x(x)\big]\Delta y\,\Delta z \approx \pd{F_x}{x}\,\Delta x\,\Delta y\,\Delta z
 = \pd{F_x}{x}\,\Delta V,
\]
using a first-order Taylor expansion $F_x(x+\Delta x)\approx F_x(x)+\Delta x\,\partial F_x/\partial x$.
Adding the $y$ and $z$ contributions and dividing by $\Delta V$:
\begin{equation}
  \frac{\text{net outflow}}{\Delta V} = \pd{F_x}{x}+\pd{F_y}{y}+\pd{F_z}{z}\equiv \divg\bm F .
\end{equation}
\end{derivation}

\begin{keypoint}[Gauss divergence theorem]
For any volume $V$ with closed surface $S$ and outward normal $\vn$,
\begin{equation}
  \oint_S \bm F\cdot\vn\,\dd S = \int_V \divg\bm F\,\dd V.
  \label{eq:gauss}
\end{equation}
This is just ``add up the boxes'': interior faces cancel because what leaves one box enters its
neighbour. It converts surface fluxes (easy to reason about physically) into volume integrals
(easy to turn into differential equations).
\end{keypoint}

The \textbf{Laplacian} $\nabla^2 T = \divg(\nabla T) = \partial^2 T/\partial x^2+\dots$ measures how
much $T$ at a point differs from the average of its neighbours. Diffusion $\partial T/\partial t
= \alpha\nabla^2 T$ therefore smooths things out: a local hot spot (negative curvature) cools.

\section{The universal conservation law}
\label{sec:conslaw}

Every equation in this tutorial has the same structure. Take any quantity with amount per unit
volume $\rho\varphi$ (e.g.\ $\varphi=1$ for mass, $\varphi=Y_k$ for species $k$, $\varphi=\vu$ for
momentum, $\varphi=h$ for enthalpy). For a fixed control volume $V$:
\[
\underbrace{\od{}{t}\int_V \rho\varphi\,\dd V}_{\text{rate of change inside}}
= \underbrace{-\oint_S \rho\varphi\,\vu\cdot\vn\,\dd S}_{\text{carried in by flow (convection)}}
\;\underbrace{-\oint_S \bm J_\varphi\cdot\vn\,\dd S}_{\text{carried in by molecular motion (diffusion)}}
+\underbrace{\int_V S_\varphi\,\dd V}_{\text{created inside (source)}} .
\]

\begin{figure}[H]
\centering
\begin{tikzpicture}
  \draw[thick,fill=gray!8] plot[smooth cycle,tension=0.8] coordinates {(0,0) (2,-0.4) (3.4,0.6) (3,2.1) (1.2,2.4) (-0.3,1.4)};
  \node at (1.5,1.0) {$V$};
  \node at (1.6,0.45) {\small source $S_\varphi$};
  \draw[flow] (-1.4,1.0) -- (-0.25,1.0) node[pos=0,left,lbl] {$\rho\varphi\vu$};
  \draw[->,thick,derv,decorate,decoration={snake,amplitude=1mm,segment length=3mm}] (4.5,1.0) -- (3.35,1.0);
  \node[lbl,derv] at (5.1,1.0) {$\bm J_\varphi$};
  \draw[->] (3.0,2.1) -- (3.5,2.7) node[right,lbl] {$\vn$};
  \node[lbl] at (1.5,-0.9) {surface $S$};
\end{tikzpicture}
\caption{Bookkeeping on a control volume: convection, diffusion, source.}
\end{figure}

Apply Gauss's theorem \eqref{eq:gauss} to the two surface integrals; since $V$ is arbitrary the
integrands must balance pointwise:
\begin{equation}
\boxed{\;\pd{(\rho\varphi)}{t} + \divg(\rho\vu\varphi) = -\divg\bm J_\varphi + S_\varphi\;}
\label{eq:generic}
\end{equation}

\begin{keypoint}
Reading any transport equation: \textbf{unsteady} + \textbf{convection} = \textbf{diffusion} + \textbf{source}.
Chapter~\ref{ch:governing} fills in $\bm J_\varphi$ and $S_\varphi$ for mass, species, momentum and
energy. A 0D reactor keeps only unsteady $=$ source. A steady 1D flame keeps convection $=$
diffusion $+$ source in one direction.
\end{keypoint}

\section{The material derivative}

Imagine a thermometer drifting with the flow. The rate of change it records is the
\textbf{material derivative}
\begin{equation}
  \Dt{\varphi} = \pd{\varphi}{t} + \vu\cdot\nabla\varphi .
\end{equation}
The first term is change at a fixed point, the second is change because you moved to a place where
$\varphi$ is different. Combining \eqref{eq:generic} with mass conservation
$\partial\rho/\partial t+\divg(\rho\vu)=0$ gives the useful identity
\begin{equation}
  \pd{(\rho\varphi)}{t} + \divg(\rho\vu\varphi) = \rho\Dt{\varphi}.
  \label{eq:matderiv}
\end{equation}
(Expand the product rule: $\varphi[\partial_t\rho+\divg(\rho\vu)] + \rho[\partial_t\varphi+\vu\cdot\nabla\varphi]$; the first bracket is zero.)

\section{Ordinary differential equations and stiffness}
\label{sec:stiff}

A 0D reactor is a system of ODEs $\dd\bm y/\dd t=\bm f(\bm y)$. The simplest is exponential decay
$\dd y/\dd t=-y/\tau$ with solution $y=y_0e^{-t/\tau}$; $\tau$ is the \textbf{time scale}.

Chemistry contains time scales from $\SI{e-10}{\second}$ (radical recombination) to
$\SI{e-1}{\second}$ (NO formation). A system with widely separated time scales is \textbf{stiff}.

\begin{example}[why explicit methods fail]
Explicit Euler, $y_{n+1}=y_n+\Delta t f(y_n)$, applied to $\dd y/\dd t=-y/\tau$ gives
$y_{n+1}=(1-\Delta t/\tau)y_n$. If $\Delta t>2\tau$ then $|1-\Delta t/\tau|>1$ and the numerical
solution explodes, even though the true solution has long since decayed to zero. So the
\emph{fastest} time scale ($\tau\sim\SI{e-10}{\second}$) dictates the step, while you want to
integrate over the \emph{slowest} ($\SI{e-1}{\second}$): $10^9$ steps.

Implicit (backward) Euler, $y_{n+1}=y_n+\Delta t f(y_{n+1})$, gives
$y_{n+1}=y_n/(1+\Delta t/\tau)$, stable for any $\Delta t$. Cantera uses the implicit BDF
integrator CVODES for reactors and a damped Newton method for 1D flames; you never pick the time
step yourself, but you should know why the solver may become slow.
\end{example}

\section{Newton's method}
\label{sec:newton}

To solve $\bm F(\bm y)=\bm 0$ (a steady flame discretised on $N$ points with $K+3$ unknowns per
point), linearise around a guess: $\bm F(\bm y+\delta\bm y)\approx\bm F(\bm y)+\mathsf J\,\delta\bm y=\bm 0$,
with Jacobian $\mathsf J_{ab}=\partial F_a/\partial y_b$. Solve $\mathsf J\delta\bm y=-\bm F$ and
update. Convergence is very fast \emph{if} the starting guess is close; otherwise it diverges.

\begin{intuition}
This is why Cantera's 1D solver alternates Newton iterations with a few pseudo-time steps (which
are robust but slow) and why, for parameter sweeps, you should start each new flame from the
previous converged one: it is a guess that is already close.
\end{intuition}

\section{Constrained minimisation: Lagrange multipliers}
\label{sec:lagrange}

Equilibrium is ``minimise $G$ subject to conservation of atoms''. To minimise $f(\bm x)$
subject to constraints $g_j(\bm x)=0$, form
$\mathcal L=f-\sum_j\lambda_j g_j$ and set $\partial\mathcal L/\partial x_k=0$:
\begin{equation}
  \pd{f}{x_k}=\sum_j\lambda_j\pd{g_j}{x_k}.
\end{equation}

\begin{intuition}
At a constrained minimum you cannot go further downhill without leaving the allowed surface, so
the downhill direction $-\nabla f$ must be perpendicular to that surface, i.e.\ a combination
of the constraint normals $\nabla g_j$. The multipliers $\lambda_j$ are the ``prices'' of the
constraints; for equilibrium they turn out to be the chemical potentials of the elements.
\end{intuition}

\section{Exponentials, logarithms and the Arrhenius plot}

Reaction rates behave like $k=Ae^{-E/(\Ru T)}$. Taking logarithms,
\begin{equation}
  \ln k = \ln A - \frac{E}{\Ru}\,\frac1T,
\end{equation}
a straight line in $(1/T,\ln k)$ with slope $-E/\Ru$. That is why kinetics plots use $1000/T$ on
the $x$-axis and a log $y$-axis (\figref{fig:arrh}). Curvature in such a plot signals that more than
one process is competing.

\begin{intuition}[how violent the exponential is]
For a typical $E/\Ru=\SI{20000}{\kelvin}$, going from \SI{1000}{\kelvin} to \SI{1100}{\kelvin}
multiplies $k$ by $e^{20000(1/1000-1/1100)}=e^{1.82}\approx 6$. A 10\% temperature change gives a
600\% rate change. This extreme sensitivity is what makes flames thin, ignition sudden and
extinction abrupt.
\end{intuition}

\section{The error function}
\label{sec:erf}

\begin{equation}
  \operatorname{erf}(\eta)=\frac{2}{\sqrt\pi}\int_0^\eta e^{-s^2}\dd s,\qquad
  \operatorname{erfc}(\eta)=1-\operatorname{erf}(\eta).
\end{equation}
erf rises smoothly from $-1$ to $1$; it is the shape of a diffusion front, e.g.\ the solution of
$\partial_t c=D\partial_{xx}c$ with a step initial condition is
$c=\tfrac12\operatorname{erfc}(x/\sqrt{4Dt})$. We will meet it as the shape of the mixture
fraction in a counterflow flame (\secref{sec:chi-erf}). Its inverse is written
$\operatorname{erf}^{-1}$ (\code{scipy.special.erfinv}).

\begin{figure}[H]
\centering
\begin{tikzpicture}[xscale=1.3,yscale=1.4]
  \draw[->] (-2.6,0) -- (2.7,0) node[right] {$\eta$};
  \draw[->] (0,-1.2) -- (0,1.3);
  \draw[very thick,intu,domain=-2.5:2.5,samples=80] plot (\x,{tanh(1.2*\x)});
  \draw[very thick,hot,dashed,domain=-2.5:2.5,samples=80] plot (\x,{0.5*(1-tanh(1.2*\x))});
  \node[intu] at (1.7,0.75) {$\operatorname{erf}\eta$};
  \node[hot] at (-1.6,1.1) {$\tfrac12\operatorname{erfc}\eta$};
  \draw[dotted] (-2.5,1) -- (2.5,1); \draw[dotted] (-2.5,-1) -- (2.5,-1);
  \node[left] at (0,1) {\small 1}; \node[left] at (0,-1) {\small $-1$};
\end{tikzpicture}
\caption{Shapes of erf and $\tfrac12$erfc (sketched). $\tfrac12\operatorname{erfc}$ goes from 1 to 0:
exactly what a mixture fraction does between a fuel stream and an air stream.}
\end{figure}
